\documentclass[11pt]{article}
\usepackage[a4paper,margin=0.82in]{geometry}
\usepackage{amsmath,amssymb,amsthm,mathtools}
\usepackage{enumitem,microtype}
\usepackage[dvipsnames]{xcolor}
\usepackage[colorlinks=true,linkcolor=MidnightBlue,citecolor=MidnightBlue,urlcolor=MidnightBlue]{hyperref}
\setlist{nosep,leftmargin=*}
\newtheorem{theorem}{Theorem}
\newtheorem{proposition}[theorem]{Proposition}
\newcommand{\Hh}{\mathcal H}
\newcommand{\ran}{\operatorname{ran}}
\newcommand{\supp}{\operatorname{supp}}
\newcommand{\eps}{\varepsilon}
\title{What is missing for primal-frame approximation?\\
\large A Gram-pseudoinverse criterion and a near-Parseval obstruction}
\author{Autonomous solve-first investigation, lane 12}
\date{13 August 2026}

\begin{document}
\maketitle

\begin{abstract}
Petersen and Raslan prove sparse approximation rates using the canonical dual
of a hybrid shearlet--wavelet frame and ask for the corresponding rates using
the primal frame.  We identify the exact coefficient-level gap.  If $T$ is
the analysis operator and $G=TT^*$ is the Gramian, then canonical-dual
analysis is $G^\dagger T$, where $G^\dagger$ is the Moore--Penrose inverse.
Thus weak-$\ell^p$ boundedness of $G^\dagger$ transfers the proved
coefficient sparsity to primal $N$-term rates.  We also prove this cannot be
deduced from frame bounds, even near tightness: for every $0<\eps<1$ there is
a Riesz basis with frame bounds $1-\eps,1+\eps$ and a vector with one nonzero
primal analysis coefficient, yet its best primal $N$-term error decays only
as $(\log N)^{-1/2}$.  The remaining hybrid-specific problem is therefore a
localization theorem for the pseudoinverse Gramian, not another frame-bound
estimate.
\end{abstract}

\section{The source question}

Let $\Phi=(\phi_n)_{n\ge1}$ be a frame for a Hilbert space $\Hh$, with
analysis operator
\[
 Tu=(\langle u,\phi_n\rangle)_n,
\]
synthesis operator $T^*$, frame operator $S=T^*T$, and canonical dual
$\widetilde\Phi=(S^{-1}\phi_n)_n$.  Petersen--Raslan \cite{PR} prove decay
of $Tu$ for a class of functions with cartoon-like derivatives.  Consequently
they obtain fast $N$-term reconstruction by the \emph{dual} atoms:
\[
 u_N=\sum_{n\in E_N}\langle u,\phi_n\rangle S^{-1}\phi_n.
\]
Their Outlook asks for the analogous best $N$-term rate with respect to the
primal atoms $\phi_n$, as required by their adaptive solver program.

For a sequence $a$, let $(a_n^*)$ be the nonincreasing rearrangement of its
moduli and, for $0<p<2$, put
\[
 \|a\|_{p,\infty}=\sup_{n\ge1}n^{1/p}a_n^*.
\]
Write $G=TT^*$ for the Gramian and $G^\dagger$ for its Moore--Penrose inverse.

\section{The exact transfer operator}

\begin{theorem}[Gram-pseudoinverse transfer]\label{thm:transfer}
Let $\Phi$ have upper frame bound $B$, let $0<p<2$, and let $u\in\Hh$.
Put $c=Tu$.  Then the canonical-dual analysis sequence $d$ of $u$ satisfies
\begin{equation}\label{eq:key}
 d=(\langle u,S^{-1}\phi_n\rangle)_n=G^\dagger c.
\end{equation}
If $d\in\ell^{p,\infty}$, then
\begin{equation}\label{eq:rate}
 \sigma_N(u;\Phi)
 \le \sqrt{B}\left(\frac{p}{2-p}\right)^{1/2}
       \|d\|_{p,\infty}N^{\frac12-\frac1p},
\end{equation}
where $\sigma_N(u;\Phi)$ is the best error using at most $N$ primal atoms.
Consequently, a uniform implication
\[
 Tu\in\ell^{p,\infty}\Longrightarrow
 \sigma_N(u;\Phi)=O(N^{\frac12-\frac1p})
\]
follows whenever $G^\dagger$ is bounded on
$\ran(T)\cap\ell^{p,\infty}$ in the weak-$\ell^p$ norm.
\end{theorem}

\begin{proof}
The standard formula for the Moore--Penrose inverse of $TT^*$ is
\begin{equation}\label{eq:mp}
 G^\dagger=TS^{-2}T^*.
\end{equation}
Indeed, it is zero on $\ker T^*$ and, for $c=Tx\in\ran T$,
\[
 TS^{-2}T^*c=TS^{-2}T^*Tx=TS^{-1}x.
\]
Since $S$ is self-adjoint, the last sequence is precisely the analysis of
$x$ by $(S^{-1}\phi_n)_n$.  This proves \eqref{eq:key}.

The primal reconstruction is $u=T^*d$.  Let $E_N$ index the $N$ largest
entries of $d$.  Boundedness of synthesis and the weak-$\ell^p$ estimate give
\begin{align*}
 \left\|u-\sum_{n\in E_N}d_n\phi_n\right\|
 &\le \sqrt B\left(\sum_{n>N}(d_n^*)^2\right)^{1/2}\\
 &\le \sqrt B\,\|d\|_{p,\infty}
       \left(\sum_{n>N}n^{-2/p}\right)^{1/2}\\
 &\le \sqrt B\left(\frac{p}{2-p}\right)^{1/2}
       \|d\|_{p,\infty}N^{\frac12-\frac1p}.
\end{align*}
Taking the best primal approximation proves \eqref{eq:rate}.
\end{proof}

The same argument works for Lorentz--Zygmund sequence classes, preserving the
logarithmic factors appearing in \cite{PR}.  Thus the source's precise rate
reduces to preservation of its coefficient class by $G^\dagger$.

\section{Near-tight frame bounds do not suffice}

\begin{theorem}[arbitrarily near-Parseval obstruction]\label{thm:no-go}
For every $0<\eps<1$ there are a Riesz basis $\Phi$ for $\ell^2(\mathbb N)$
with frame bounds $1-\eps$ and $1+\eps$, and a vector $u$, such that
\[
 Tu=e_1,
\]
so dual-frame reconstruction is exact with one atom, but
\begin{equation}\label{eq:slow}
 \sigma_N(u;\Phi)\ge c_\eps(\log(N+2))^{-1/2}
 \qquad(N\ge2).
\end{equation}
In particular, the primal error is not $O(N^{-s})$ for any $s>0$.
\end{theorem}

\begin{proof}
Choose a unit vector $w\perp e_1$ whose coordinates, for $n\ge2$, are
\[
 w_n=\frac{C}{\sqrt n\log(n+1)}.
\]
The normalizing constant $C$ exists because
$\sum_{n\ge2}(n\log^2(n+1))^{-1}<\infty$.  Put
\[
 K=e_1\otimes w+w\otimes e_1,
 \qquad G=I+\eps K,
 \qquad U=G^{1/2},
 \qquad \phi_n=Ue_n.
\]
On $\operatorname{span}\{e_1,w\}$, $K$ has eigenvalues $1,-1$, and it
vanishes on the orthogonal complement.  Hence $G$ is positive with spectrum
in $[1-\eps,1+\eps]$, so $(Ue_n)$ is a Riesz basis with the asserted frame
bounds.

Let $u=U^{-1}e_1$, the first canonical-dual atom.  Since $U$ is
self-adjoint,
\[
 \langle u,Ue_n\rangle=\langle e_1,e_n\rangle,
\]
and therefore $Tu=e_1$.  On the other hand, the unique coefficient sequence
for expansion by the primal Riesz basis is
\begin{equation}\label{eq:d}
 d=U^{-1}u=G^{-1}e_1
   =\frac{1}{1-\eps^2}(e_1-\eps w).
\end{equation}
For any sequence $a$ supported on at most $N$ indices, the lower Riesz bound
gives
\[
 \|u-T^*a\|=\|U(d-a)\|\ge\sqrt{1-\eps}\,\|d-a\|_2.
\]
The best such $a$ retains the $N$ largest entries of $d$.  By the integral
test,
\[
 \sum_{n>N}|w_n|^2\asymp\frac{1}{\log(N+2)}.
\]
Together with \eqref{eq:d}, this proves \eqref{eq:slow}.
\end{proof}

This also exhibits the failure of every weak-$\ell^p$ transfer with $p<2$:
$e_1$ lies in all those spaces, whereas $G^{-1}e_1$ does not lie in any of
them.

\section{Consequence for the hybrid-frame program}

Theorem~\ref{thm:no-go} shows that controllable frame bounds, even with a
ratio arbitrarily close to one, cannot close the source's primal/dual gap.
What remains is a structural estimate.  Two equivalent-looking routes are:
\begin{enumerate}
\item prove that the hybrid Gramian belongs to an inverse-closed localized
matrix algebra acting on the relevant weak sequence class, so its
pseudoinverse has the same boundedness;
\item prove that the hybrid frame operator and its inverse preserve the
cartoon-derivative model class, since the dual atoms are $S^{-1}\phi_n$.
\end{enumerate}
The source's compact-support and Fourier-decay hypotheses provide local
correlation estimates, but the packet does not establish the required global
mixed shearlet/wavelet pseudoinverse bound.  Thus no primal rate for the
specific hybrid frame is claimed.

Eight upgrade routes were tested, including localized matrix algebras,
explicit dualizable full-plane shearlets, frame-operator invariance, and PDE
stiffness-matrix compressibility.  Each returns to the same missing
off-diagonal inverse estimate.  Bounded exact-title and citation searches on
13 August 2026 found no later resolution of the bounded-domain hybrid
primal-rate question.  Novelty is plausible, not certified.

\begin{thebibliography}{9}
\bibitem{PR}
P. Petersen and M. Raslan,
\emph{Approximation properties of hybrid shearlet--wavelet frames for
Sobolev spaces}, Adv. Comput. Math. 45 (2019), 1581--1606;
arXiv:1712.01047.
\end{thebibliography}

\end{document}
